% Source: https://artofproblemsolving.com/community/c3h1722841_modulo_problem
% Stable identity: aops-topic-1722841-q01
\textbf{The six squares are distinct roots of $x^6-1$.} Modulo $13$ the squares are $1,4,9,16\equiv3,25\equiv12,36\equiv10$: six different residues $r_1,\dots,r_6$. Each is $r=j^2$ with $13\nmid j$, so by Fermat's little theorem $r^6=j^{12}\equiv1$. Thus every $r_i$ is a root of $x^6-1$ modulo $13$.

\textbf{Comparing the polynomials.} Let $P(x)=\prod_{i=1}^{6}(x-r_i)$ and $D(x)=P(x)-(x^6-1)$. Both are monic of degree $6$, so $D$ has degree at most $5$, and $D(r_i)\equiv0-0=0$ for the six distinct residues $r_i$. Over the integers modulo the prime $13$, a nonzero polynomial of degree at most $5$ has at most $5$ roots (if $D(r)\equiv0$, then $D(x)\equiv(x-r)E(x)$, and since $13$ is prime a product is $0$ only if a factor is, so roots of $D$ other than $r$ are roots of $E$; induct on the degree). Hence all coefficients of $D$ are divisible by $13$, which is the claim. $\blacksquare$

(The residues $1,3,4,9,10,12$ are the quadratic residues modulo $13$, and the identity is a case of Euler's criterion.)
